\documentclass[11pt,a4paper]{article}

% --- PAQUETES DE IDIOMA Y CODIFICACIÓN ---
\usepackage[utf8]{inputenc}
\usepackage[T1]{fontenc}
\usepackage[spanish,es-tabla,es-nodecimaldot,es-noshorthands]{babel}

% --- MATEMÁTICAS Y SÍMBOLOS ---
\usepackage{amsmath,amssymb,amsfonts}
\usepackage{mathtools}

% --- FORMATO Y DISEÑO DE PÁGINA ---
\usepackage[top=2.5cm, bottom=2.5cm, left=2.5cm, right=2.5cm]{geometry}
\usepackage{fancyhdr}
\usepackage{booktabs}
\usepackage{array}
\usepackage{enumitem}
\usepackage{titlesec}

% --- COLORES Y CÓDIGO ---
\usepackage[table]{xcolor}
\definecolor{pujblue}{RGB}{0, 51, 102}
\definecolor{darkblue}{RGB}{16, 44, 87}
\definecolor{codebg}{RGB}{248, 249, 250}
\definecolor{darkgray}{RGB}{60, 60, 60}

\usepackage{listings}
\usepackage[colorlinks=true, linkcolor=pujblue, citecolor=pujblue, urlcolor=pujblue]{hyperref}

% --- CONFIGURACIÓN DE LISTINGS (SOPORTE COMPLETO UTF-8 / ESPAÑOL) ---
\lstset{
    language=Python,
    backgroundcolor=\color{codebg},
    basicstyle=\ttfamily\footnotesize,
    keywordstyle=\color{pujblue}\bfseries,
    commentstyle=\color{gray}\itshape,
    stringstyle=\color{teal},
    breaklines=true,
    frame=single,
    rulecolor=\color{black!20},
    numbers=left,
    numberstyle=\tiny\color{gray},
    numbersep=8pt,
    tabsize=4,
    showstringspaces=false,
    literate=
      {á}{{\'a}}1 {é}{{\'e}}1 {í}{{\'i}}1 {ó}{{\'o}}1 {ú}{{\'u}}1
      {Á}{{\'A}}1 {É}{{\'E}}1 {Í}{{\'I}}1 {Ó}{{\'O}}1 {Ú}{{\'U}}1
      {ñ}{{\~n}}1 {Ñ}{{\~N}}1 {ü}{{\"u}}1 {Ü}{{\"U}}1
      {¿}{{?`}}1 {¡}{{!`}}1
}

% --- ENCABEZADOS Y PIES DE PÁGINA ---
\pagestyle{fancy}
\fancyhf{}
\fancyhead[L]{\small\color{darkgray}\textit{Pontificia Universidad Javeriana Cali --- Computación Científica}}
\fancyhead[R]{\small\color{darkgray}\textit{Taller 2: Factorización PLU}}
\fancyfoot[C]{\thepage}
\renewcommand{\headrulewidth}{0.4pt}
\renewcommand{\footrulewidth}{0.4pt}

\titleformat{\section}{\Large\bfseries\color{pujblue}}{\thesection}{1em}{}[\titlerule]
\titleformat{\subsection}{\large\bfseries\color{darkblue}}{\thesubsection}{1em}{}
\titleformat{\subsubsection}{\normalsize\bfseries\color{darkblue}}{\thesubsubsection}{1em}{}

% ==============================================================================
% INICIO DEL DOCUMENTO
% ==============================================================================
\begin{document}

\begin{center}
    {\LARGE\bfseries\color{pujblue} Taller 2: Modelo Económico Multisectorial mediante Factorización PLU}\\[0.3cm]
    {\large Computación Científica --- Pontificia Universidad Javeriana Cali}\\[0.1cm]
    {\normalsize Profesor: Juan Fernando Paz}
\end{center}
\vspace{0.3cm}

% ==============================================================================
% SECCIÓN 1: MODELO Y ALGORITMOS EN PYTHON
% ==============================================================================
\section{Implementación de la Factorización PLU y Sustituciones en Python}

Se analiza una economía simplificada de $30$ sectores productivos mediante el sistema lineal $Ax = b$, donde $A \in \mathbb{R}^{30 \times 30}$ es la matriz de requerimientos entre sectores, $x \in \mathbb{R}^{30}$ es el nivel de producción de cada sector y $b \in \mathbb{R}^{30}$ es el vector de demanda final.

A continuación se presenta la implementación de la factorización $PA = LU$ mediante eliminación de Gauss con pivoteo parcial y los algoritmos de sustitución hacia adelante ($Ly = Pb$) y hacia atrás ($Ux = y$):

\begin{lstlisting}
import time
import numpy as np
from scipy.linalg import lu, solve_triangular

np.set_printoptions(precision=4, suppress=True)

def factorizacion_plu(A):
    """
    Calcula la descomposicion PA = LU mediante eliminacion de Gauss
    con pivoteo parcial.
    """
    A = np.array(A, dtype=float)
    n = len(A)

    U = np.copy(A)
    L = np.eye(n, dtype=float)
    P = np.eye(n, dtype=float)

    for k in range(n - 1):
        # 1. Buscar el pivote de mayor valor absoluto en la columna k (desde fila k hasta n-1)
        piv = k
        mayor = abs(U[k, k])
        for i in range(k + 1, n):
            if abs(U[i, k]) > mayor:
                mayor = abs(U[i, k])
                piv = i

        # 2. Si la fila con mayor valor no es la fila k, intercambiar filas
        if piv != k:
            temp_U = np.copy(U[k])
            U[k] = U[piv]
            U[piv] = temp_U

            temp_P = np.copy(P[k])
            P[k] = P[piv]
            P[piv] = temp_P

            for col in range(k):
                temp_L = L[k, col]
                L[k, col] = L[piv, col]
                L[piv, col] = temp_L

        # 3. Eliminacion hacia adelante debajo del pivote U[k, k]
        for i in range(k + 1, n):
            factor = U[i, k] / U[k, k]
            L[i, k] = factor
            for j in range(k, n):
                U[i, j] = U[i, j] - factor * U[k, j]

    return P, L, U

def sustitucion_adelante(L, b):
    L = np.array(L, dtype=float)
    b = np.array(b, dtype=float)
    n = len(b)
    y = np.zeros(n, dtype=float)

    for i in range(n):
        suma = 0.0
        for k in range(i):
            suma += L[i, k] * y[k]
        y[i] = (b[i] - suma) / L[i, i]

    return y

def sustitucion_atras(U, y):
    U = np.array(U, dtype=float)
    y = np.array(y, dtype=float)
    n = len(y)
    x = np.zeros(n, dtype=float)

    for i in range(n - 1, -1, -1):
        suma = 0.0
        for k in range(i + 1, n):
            suma += U[i, k] * x[k]
        x[i] = (y[i] - suma) / U[i, i]

    return x

def resolver_sistema_plu(P, L, U, b):
    Pb = P @ b
    y = sustitucion_adelante(L, Pb)
    x = sustitucion_atras(U, y)
    return x, y, Pb
\end{lstlisting}

% ==============================================================================
% SECCIÓN 2: ACTIVIDADES (1 A 6)
% ==============================================================================
\section{Desarrollo de las Actividades}

\subsection{Actividad 1: ¿Por qué es necesario utilizar pivoteo parcial?}
Al inspeccionar la diagonal principal de la matriz de relaciones $A \in \mathbb{R}^{30 \times 30}$, se observa que su primer elemento es exactamente cero ($a_{0,0} = 0.0$), y además existen ceros en la diagonal en las filas con índices:
\[
    \text{Filas con } a_{ii} = 0: \quad [0, \; 5, \; 10, \; 15, \; 20, \; 25]
\]
\begin{enumerate}[leftmargin=*]
    \item \textbf{Imposibilidad de aplicar LU directo (división por cero):} En la eliminación de Gauss estándar o en el método de Crout sin pivoteo, el primer paso divide entre el pivote $a_{0,0} = 0$, lo que ocasiona una división por cero inmediata aunque la matriz $A$ sea invertible.
    \item \textbf{Estabilidad numérica:} El pivoteo parcial busca en cada etapa $k$ el elemento de mayor valor absoluto en la columna actual desde la fila $k$ hacia abajo ($\max_{i \ge k} |u_{ik}|$) e intercambia las filas mediante la matriz de permutación $P$. Esto garantiza que todos los multiplicadores en $L$ cumplan $|l_{ik}| \le 1$ y evita la amplificación de errores de redondeo.
\end{enumerate}

\subsection{Actividad 2: Obtención de las Matrices $P$, $L$ y $U$ ($PA = LU$)}
Al ejecutar \texttt{P, L, U = factorizacion\_plu(A)}, se obtienen tres matrices de dimensión $30 \times 30$. A continuación se muestra el bloque superior izquierdo ($5 \times 5$) de $L$ y $U$:
\begin{equation}
    L_{0:5, 0:5} = \begin{pmatrix}
        1.0000 & 0.0000 & 0.0000 & 0.0000 & 0.0000 \\
        0.2222 & 1.0000 & 0.0000 & 0.0000 & 0.0000 \\
        0.6667 & -0.1304 & 1.0000 & 0.0000 & 0.0000 \\
        0.7778 & -0.4783 & 0.3621 & 1.0000 & 0.0000 \\
        0.5556 & -0.1739 & 0.9368 & 0.0653 & 1.0000
    \end{pmatrix}
\end{equation}
\begin{equation}
    U_{0:5, 0:5} = \begin{pmatrix}
        9.0000 & 6.0000 & 3.0000 & 1.0000 & 4.0000 \\
        0.0000 & 7.6667 & 4.3333 & 5.7778 & 3.1111 \\
        0.0000 & 0.0000 & 7.5652 & 4.0870 & -1.2609 \\
        0.0000 & 0.0000 & 0.0000 & 9.5057 & -0.1667 \\
        0.0000 & 0.0000 & 0.0000 & 0.0000 & 8.5109
    \end{pmatrix}
\end{equation}

\subsection{Actividad 3: Verificación Numérica de la Factorización $\|PA - LU\|_2$}
Calculando la norma euclidiana de la diferencia entre $PA$ y $LU$:
\begin{equation}
    \|P A - L U\|_2 = 1.2223 \times 10^{-14}
\end{equation}
Dado que el valor es del orden de $10^{-14}$ (precisión de punto flotante de 64 bits), se verifica que se cumple $PA = LU$ con exactitud numérica.

\subsection{Actividades 4, 5 y 6: Resolución de los Escenarios Económicos y Residuos}
Se generaron tres vectores de demanda $b^{(1)}, b^{(2)}, b^{(3)} \in \mathbb{R}^{30}$ dentro del intervalo $[50, 500]$:
\begin{lstlisting}
np.random.seed(42)
# Escenario 1 (Normal): valores en [180, 320]
b1_normal = np.random.uniform(180.0, 320.0, size=30)
# Escenario 2 (Expansion): aumento de demanda respecto al normal, dentro de [240, 500]
b2_expansion = b1_normal + np.random.uniform(60.0, 180.0, size=30)
# Escenario 3 (Desaceleracion): disminucion de demanda respecto al normal, dentro de [50, 260]
b3_desaceleracion = b1_normal - np.random.uniform(60.0, 130.0, size=30)
\end{lstlisting}

Para cada escenario $k \in \{1, 2, 3\}$ se transformó el sistema $Ax^{(k)} = b^{(k)}$ en $LUx^{(k)} = Pb^{(k)}$, resolviendo primero $Ly^{(k)} = Pb^{(k)}$ (sustitución hacia adelante) y luego $Ux^{(k)} = y^{(k)}$ (sustitución hacia atrás), verificando cada solución con el residuo $r^{(k)} = Ax^{(k)} - b^{(k)}$:

\begin{table}[h!]
    \centering
    \small
    \caption{Resumen de resultados y verificación de residuos para los tres escenarios.}
    \begin{tabular}{lccc}
        \toprule
        \textbf{Escenario Económico} & \textbf{Rango de $b^{(k)}$} & \textbf{Muestra $x^{(k)}_{0:3}$} & \textbf{Norma del Residuo $\|r^{(k)}\|_2$} \\
        \midrule
        1. Condiciones normales & $[182.88, \, 315.79]$ & $[4.7360, \, 7.6456, \, 37.3408]$ & $9.3921 \times 10^{-13}$ \\
        2. Expansión económica & $[257.53, \, 460.55]$ & $[22.6934, \, 3.6774, \, 70.9056]$ & $1.7912 \times 10^{-12}$ \\
        3. Desaceleración económica & $[68.82, \, 241.88]$ & $[7.2237, \, 1.8560, \, 26.2147]$ & $5.0223 \times 10^{-13}$ \\
        \bottomrule
    \end{tabular}
\end{table}

En los tres escenarios, $\|r^{(k)}\|_2 \approx 10^{-13}$, comprobando que los vectores $x^{(1)}, x^{(2)}, x^{(3)}$ satisfacen el sistema lineal con precisión de máquina.

% ==============================================================================
% SECCIÓN 3: ANÁLISIS ECONÓMICO
% ==============================================================================
\section{Análisis Económico}

\begin{enumerate}[leftmargin=*]
    \item \textbf{¿Qué sectores requieren un mayor nivel de producción en cada escenario?}
    \begin{itemize}
        \item \textbf{Escenario 1 (Normal):} Sector 12 (índice 11) con $44.6787$, Sector 15 (índice 14) con $41.3711$ y Sector 30 (índice 29) con $38.8549$ unidades.
        \item \textbf{Escenario 2 (Expansión):} Sector 12 (índice 11) con $95.6412$, Sector 9 (índice 8) con $87.4098$ y Sector 30 (índice 29) con $81.9160$ unidades.
        \item \textbf{Escenario 3 (Desaceleración):} Sector 12 (índice 11) con $42.7345$, Sector 15 (índice 14) con $31.8176$ y Sector 30 (índice 29) con $26.4769$ unidades.
    \end{itemize}
    Los sectores 12 y 30 encabezan de manera consistente los requerimientos productivos al actuar como sectores proveedores transversales de insumos para el resto de la economía.

    \item \textbf{¿Qué sectores presentan el mayor cambio en su nivel de producción entre el escenario normal y el escenario de expansión?}
    \begin{itemize}
        \item \textbf{Sector 11 (índice 10):} variación de $-52.6194$ unidades (magnitud $52.6194$).
        \item \textbf{Sector 9 (índice 8):} incremento de $+52.4702$ unidades (magnitud $52.4702$).
        \item \textbf{Sector 12 (índice 11):} incremento de $+50.9625$ unidades (magnitud $50.9625$).
    \end{itemize}

    \item \textbf{¿Qué sectores se ven más afectados durante la desaceleración económica?}
    \begin{itemize}
        \item \textbf{Sector 4 (índice 3):} variación de $+18.8072$ unidades (magnitud $18.8072$).
        \item \textbf{Sector 30 (índice 29):} caída de $-12.3780$ unidades (magnitud $12.3780$).
        \item \textbf{Sector 13 (índice 12):} caída de $-11.9702$ unidades (magnitud $11.9702$).
    \end{itemize}

    \item \textbf{¿Por qué la matriz $A$ puede mantenerse constante mientras cambia el vector $b$?}
    Porque la matriz $A$ representa las relaciones tecnológicas y productivas internas entre los 30 sectores (cuántas unidades de insumo requiere cada sector de los demás por unidad producida). Estas proporciones dependen de la maquinaria, infraestructura y procesos industriales instalados, que permanecen fijos en el corto y mediano plazo. Por el contrario, el vector $b$ representa la demanda final del mercado (consumo, inversión y exportaciones), la cual varía continuamente según la coyuntura económica.

    \item \textbf{Si durante el día se reciben 100 nuevos vectores de demanda, ¿sería necesario volver a calcular las matrices $P$, $L$ y $U$?}
    \textbf{No.} Como la matriz de coeficientes $A$ no cambia, la descomposición $PA = LU$ depende exclusivamente de $A$ y se calcula \textbf{una única vez}. Las matrices $P$, $L$ y $U$ se almacenan en memoria y se reutilizan para los 100 vectores.

    \item \textbf{Explique qué operaciones deben realizarse nuevamente para cada nuevo vector $b$.}
    Para cada nuevo vector de demanda $b$ solo se realizan tres pasos de costo $\mathcal{O}(n^2)$:
    \begin{enumerate}
        \item Multiplicar la matriz de permutación por el vector: $\tilde{b} = P b$ (reordenar las entradas de $b$ según los intercambios de filas registrados en $P$).
        \item Resolver el sistema triangular inferior $Ly = \tilde{b}$ mediante \textbf{sustitución hacia adelante}.
        \item Resolver el sistema triangular superior $Ux = y$ mediante \textbf{sustitución hacia atrás}.
    \end{enumerate}
\end{enumerate}

% ==============================================================================
% SECCIÓN 4: ANÁLISIS DE TIEMPO
% ==============================================================================
\section{Análisis de Tiempo: Matriz Inversa vs. Factorización PLU}

Se comparó el tiempo de ejecución de calcular la solución analítica $x = A^{-1}b$ usando \texttt{Ainv = np.linalg.inv(A)} frente a la solución mediante factorización PLU con \texttt{from scipy.linalg import lu}:

\begin{lstlisting}
# 1. Solucion con Matriz Inversa
inicio_inv = time.perf_counter()
Ainv = np.linalg.inv(A)
x_inv = Ainv @ b_prueba
tiempo_inv = time.perf_counter() - inicio_inv

# 2. Solucion con Factorizacion PLU (scipy.linalg.lu)
inicio_lu = time.perf_counter()
P_scipy, L_scipy, U_scipy = lu(A)
Pb_scipy = P_scipy.T @ b_prueba
y_scipy = solve_triangular(L_scipy, Pb_scipy, lower=True, unit_diagonal=True)
x_lu = solve_triangular(U_scipy, y_scipy, lower=False)
tiempo_lu = time.perf_counter() - inicio_lu
\end{lstlisting}

\begin{table}[h!]
    \centering
    \caption{Comparación de tiempos de ejecución entre \texttt{np.linalg.inv(A)} y \texttt{scipy.linalg.lu(A)}.}
    \begin{tabular}{lcc}
        \toprule
        \textbf{Caso Evaluado} & \textbf{Con Inversa (\texttt{np.linalg.inv})} & \textbf{Con Factorización PLU (\texttt{scipy.linalg.lu})} \\
        \midrule
        Resolución de 1 escenario & $0.000031$ s & $0.000069$ s \\
        Resolución de 100 escenarios de demanda & $0.001471$ s & $0.001205$ s \\
        Diferencia numérica $\|x_{\text{inv}} - x_{\text{lu}}\|_2$ & \multicolumn{2}{c}{$3.4982 \times 10^{-13}$} \\
        \bottomrule
    \end{tabular}
\end{table}

\noindent\textbf{Conclusión del Análisis de Tiempo:}
\begin{itemize}[leftmargin=*]
    \item \textbf{Complejidad algorítmica:} Calcular la matriz inversa $A^{-1}$ requiere $\approx 2n^3$ operaciones en punto flotante, lo cual es \textbf{tres veces más costoso} que calcular la factorización $PA = LU$ ($\approx \frac{2}{3}n^3$ operaciones).
    \item \textbf{Escenarios múltiples:} Al recibir múltiples vectores de demanda $b$ (como los 100 vectores diarios), la factorización PLU descompone $A$ una sola vez y resuelve cada nuevo vector en $\mathcal{O}(2n^2)$ mediante sustituciones triangulares, superando en velocidad al recálculo con la inversa ($0.001205\text{ s}$ frente a $0.001471\text{ s}$).
    \item \textbf{Estabilidad numérica:} En computación científica se evita invertir matrices explícitamente porque $A^{-1}$ introduce mayor acumulación de errores de redondeo que resolver directamente los sistemas triangulares $Ly = Pb$ y $Ux = y$.
\end{itemize}

\end{document}
